\section{Cadre et notations}\label{sec-cadre}

\begin{objectif}
Fixer les notations et établir les trois outils dont la suite a besoin : la formule du forward, l'écart de portage,
et le passage des données Totem au repo implicite.
\end{objectif}

\subsection{Conventions}
Le temps est mesuré en années à partir de la date de valorisation, et les taux sont des taux zéro-coupon en
composition continue. Les scalaires et les vecteurs sont notés en minuscules, les matrices en majuscules ;
$\tr{A}$ désigne la transposée de $A$, $\M_{N,m}(\R)$ l'espace des matrices à $N$ lignes et $m$ colonnes,
$\sprod{x}{y}=\tr{x}\,y$ le produit scalaire de $\R^N$ et $\norm{x}$ la norme associée. L'exposant $\tot$ désigne une
grandeur issue de Totem et l'exposant $\fo$ une grandeur du front office. Le point de base (pb) vaut $10^{-4}$ ; un
écart de portage s'exprime en points de base du spot. Les notations sont récapitulées en \cref{app-notation}.

\subsection{Forward et dividendes discrets}
On note $S_0$ le spot, $r(T)$ le taux zéro-coupon de maturité $T$, et $q(T)$ le taux zéro-coupon de \emph{repo},
c'est-à-dire le coût d'emprunt du titre : détenir l'action financée par un repo rapporte $r-q$ par unité de temps.
Les dividendes sont discrets : un montant $d_i$ est détaché à la date $t_i$. On note
$\tilde d_i=d_i\,e^{-(r(t_i)-q(t_i))\,t_i}$ le dividende actualisé au taux de portage.

\begin{definition}[Forward]\label{def-forward}
Le forward de maturité $T$ vaut
\begin{equation}\label{eq-forward}
F(T)=\bigl(S_0-D(T)\bigr)\,e^{(r(T)-q(T))\,T}, \qquad D(T)=\sum_{t_i\le T}\tilde d_i .
\end{equation}
$D(T)$ est la somme des dividendes détachés avant $T$, actualisés au taux de portage.
\end{definition}

C'est l'argument de réplication habituel : on achète l'action en la finançant au taux $r-q$ et on reporte les
dividendes perçus. Le repo intervient aussi, faiblement, dans $D$ par l'actualisation des dividendes ; le bootstrap de
la \cref{sec-bootstrap} traite ce couplage exactement.

\begin{lemma}[Logarithme du forward]\label{lem-logf}
$\ln F(T)=\bigl(r(T)-q(T)\bigr)\,T+\ln\bigl(S_0-D(T)\bigr)$. Par conséquent, des perturbations du forward, des
dividendes et du repo vérifient, au premier ordre,
\begin{equation}\label{eq-additivity}
\delta\bigl(q(T)\,T\bigr)=\delta\bigl(r(T)\,T\bigr)-\delta\ln F(T)-\frac{\delta D(T)}{S_0-D(T)} .
\end{equation}
\end{lemma}
\begin{proof}
On prend le logarithme de \eqref{eq-forward}, puis on différentie, avec $\delta\ln(S_0-D)=-\delta D/(S_0-D)+O(\delta D^2)$.
\end{proof}

\subsection{L'écart de portage}
Les maturités Totem, ou \emph{piliers}, sont $T_1<\dots<T_N$ (ici $N=11$, d'un mois à cinq ans). Pour toute grandeur
$x$, on note $x_k=x(T_k)$.

\begin{definition}[Écart de portage]\label{def-delta}
L'\emph{écart de portage} au pilier $T_k$ est
\begin{equation}\label{eq-delta}
\delta_k=\bigl(q^{\tot}_k-q^{\fo}_k\bigr)\,T_k, \qquad k\in\ient1N,
\end{equation}
où $q^{\tot}$ est le repo implicite Totem (\cref{def-implied}) et $q^{\fo}$ le repo du FO. On note
$\delta=(\delta_1, \dots, \delta_N)\in\R^N$.
\end{definition}

\begin{proposition}[Le repo, différence du forward et des dividendes]\label{prop-delta}
Le contrôleur utilise la même courbe des taux des deux côtés. Alors, au couplage près par l'actualisation des
dividendes,
\begin{equation}\label{eq-delta-exact}
\delta_k=-\ln\frac{F^{\tot}_k}{F^{\fo}_k}+\ln\frac{S_0-D^{\tot}_k}{S_0-D^{\fo}_k}
\;\approx\;-\frac{F^{\tot}_k-F^{\fo}_k}{F_k}-\frac{D^{\tot}_k-D^{\fo}_k}{S_0-D_k} .
\end{equation}
\end{proposition}
\begin{proof}
Par construction, le repo implicite Totem reproduit le forward Totem avec les dividendes Totem, et le repo FO le
forward FO avec les dividendes FO. On applique le \cref{lem-logf} des deux côtés et on soustrait : les termes de taux
s'éliminent. L'approximation vient de $\ln(1+x)=x+O(x^2)$.
\end{proof}

L'écart de repo, exprimé en portage, est donc l'écart relatif de forward moins l'écart relatif de dividendes. C'est
déjà, sous une forme élémentaire, l'intuition du contrôleur : contrôle repo $=$ contrôle forward $-$ contrôle
dividendes. La \cref{sec-controle} en fait une identité sur les lignes de contrôle.

\begin{retenir}
Pourquoi travailler en portage plutôt qu'en repo : dans $\delta$, les effets des différents mécanismes
\emph{s'additionnent}, et rien n'est amplifié. Une erreur constante $\varepsilon$ sur $\ln F$ est une constante en
portage et une hyperbole $\varepsilon/T$ en repo. Tous les tests de cette note portent donc sur $\delta$.
\end{retenir}

\subsection{Les données Totem}
\begin{assumption}[Contributeurs cohérents]\label{hyp-coherent}
Le contributeur $j$ a une vue $q_j(\cdot)$ du repo et une vue $D_j(\cdot)$ des dividendes, et cote des prix d'options
cohérents avec le forward $F_j(T)=\bigl(S_0-D_j(T)\bigr)\,e^{(r(T)-q_j(T))T}$, le spot et la courbe des taux étant
communs.
\end{assumption}

\begin{assumption}[Panels]\label{hyp-panels}
Au pilier $T_k$, le panel options $\Popt_k$ compte $n^{\opt}_k$ contributeurs ; pour la période $p$, le panel
dividendes $\Pdiv_p$ en compte $n^{\dv}_p$. Les panels sont emboîtés en maturité : un contributeur qui cote une
maturité longue cote aussi les plus courtes. Les erreurs de vue sont indépendantes d'un contributeur à l'autre. Les
deux services ont $n^{\com}$ contributeurs en commun ($n^{\com}=0$ pour des panels disjoints).
\end{assumption}

\begin{assumption}[Consensus]\label{hyp-consensus}
Le consensus est la moyenne arithmétique sur le panel (après un éventuel filtre des cotations aberrantes). Totem publie,
par maturité et par période, le consensus, la dispersion entre contributeurs et leur nombre.
\end{assumption}

\subsection{Du consensus au repo implicite}\label{sec-bootstrap}
Le passage se fait en trois étapes.

\paragraph{Forwards par la parité call--put.}
\begin{lemma}[Parité et forward de consensus]\label{lem-parity}
Si chaque contributeur cote des calls et des puts européens avec un même facteur d'actualisation $B(0, T)$, alors
$C_j(K)-P_j(K)=B(0, T)\,(F_j-K)$ pour tout strike $K$. Si le consensus moyenne les mêmes contributeurs à chaque strike,
$\bar C(K)-\bar P(K)=B(0, T)\,(\bar F-K)$, où $\bar F$ est la moyenne des forwards des contributeurs. Le strike où
$\bar C=\bar P$ est donc $\bar F$, et il ne dépend pas de $B(0, T)$.
\end{lemma}
\begin{proof}
La première égalité est la parité call--put ; la deuxième s'obtient en moyennant la première, qui est affine ; la
troisième en annulant la deuxième.
\end{proof}
En pratique, on ajuste par moindres carrés pondérés la droite $K\mapsto\bar C(K)-\bar P(K)$ ; sa racine est le forward
Totem $F^{\tot}_k$ et sa pente donne un facteur d'actualisation implicite, qui sert de diagnostic. Deux situations
mettent le lemme en défaut : un filtre des cotations aberrantes appliqué séparément par strike et par type d'option (les
contributeurs moyennés ne sont alors plus les mêmes) et l'exercice américain. Un exemple chiffré et le cas de facteurs
d'actualisation différents d'un contributeur à l'autre figurent en \cref{app-parity}.

\paragraph{Allocation des dividendes.} Les niveaux Totem $\ell_p$ sont des sommes de dividendes détachés dans la
période $I_p=\intoc{e_{p-1}}{e_p}$, $p\in\ient1m$ (ici $m=6$ : la fin de 2026, puis 2027 à 2031). Le contrôleur les
répartit sur les dates de détachement de la prévision FO : un profil attribue un poids $\pi_i\ge0$ à chaque date, avec
$\sum_{t_i\in I_p}\pi_i=1$, et le dividende alloué vaut $d_i=\ell_p\,\pi_i$.

\begin{lemma}[Allocation]\label{lem-allocation}
À repo fixé,
\[
D_k=\sum_{p=1}^m a_{kp}\,\ell_p, \qquad a_{kp}=\sum_{t_i\in I_p,\ t_i\le T_k}\pi_i\,e^{-(r(t_i)-q(t_i))\,t_i} .
\]
On note $A=(a_{kp})\in\M_{N,m}(\R)$ la matrice d'allocation.
\end{lemma}
\begin{proof}
On remplace $d_i$ par $\ell_p\pi_i$ dans \eqref{eq-forward} et on regroupe les termes par période.
\end{proof}
Le coefficient $a_{kp}$ est la part, actualisée, de la période $p$ détachée avant $T_k$. Il est nul si la période
commence après $T_k$, et égal au facteur d'actualisation de la période (proche de $1$) si elle est entièrement détachée
avant $T_k$ ; entre les deux, pour une période qui chevauche un pilier, il dépend du profil : c'est précisément là que
le profil compte.

\paragraph{Bootstrap.} On cherche une courbe de repo à repo forward constant entre deux piliers, c'est-à-dire telle que
$t\mapsto q(t)\,t$ soit affine sur chaque $[T_{k-1}, T_k]$, qui reproduise les forwards. On la construit pilier par
pilier : connaissant $q_1, \dots, q_{k-1}$, l'inconnue est $x=q_kT_k$, et le forward du modèle en $T_k$ est une fonction
$\mathcal F_k(x)$, explicitée en \cref{app-bootstrap}.

\begin{proposition}[Le bootstrap est bien posé]\label{prop-bootstrap}
On suppose $S_0-D_{k-1}>0$. La fonction $\mathcal F_k$ vérifie $\mathcal F_k'\le-\mathcal F_k$ ; elle est donc
strictement décroissante là où elle est positive. Pour tout $F_k>0$, l'équation $\mathcal F_k(x)=F_k$ a une unique
solution, et la méthode de Newton appliquée à $x\mapsto\ln\mathcal F_k(x)-\ln F_k$, sécurisée par dichotomie, converge
vers elle.
\end{proposition}
La démonstration, un exemple à la main et l'algorithme sont en \cref{app-bootstrap}.

\begin{definition}[Repo implicite Totem]\label{def-implied}
Le \emph{repo implicite Totem} $q^{\tot}$ est la solution du bootstrap pour les forwards $F^{\tot}_k$ et les niveaux
de dividendes Totem $\ell^{\tot}_p$ alloués sur le profil FO. On note aussi $\hat q$ la solution pour les forwards
Totem et les dividendes \emph{FO} : elle ne dépend de Totem qu'à travers ses forwards.
\end{definition}

\begin{retenir}
\begin{itemize}[leftmargin=12pt]
\item $\delta_k=(q^{\tot}_k-q^{\fo}_k)\,T_k$ : l'écart de repo exprimé en portage, en points de base du spot.
\item $\delta_k\approx-$ (écart relatif de forward) $-$ (écart relatif de dividendes) : le repo implicite est un
résidu.
\item Le passage de Totem au repo implicite (parité, allocation, bootstrap) est exact et bien posé ; les difficultés
viennent des données, pas de l'algorithme.
\end{itemize}
\end{retenir}
